% La liste des retraits et des corrections, racontée honnêtement.
\chapter{Ce que nous avons retiré}
\markboth{Ce que nous avons retiré}{Ce que nous avons retiré}

Voici, dans l'ordre des dates, les affirmations de l'année qui ont été retirées, réfutées ou corrigées en public. Elles
ne sont pas cachées dans une note : elles font partie de l'histoire, et souvent de la meilleure partie. Chaque entrée
dit ce qui avait été affirmé, ce qui l'a fait tomber, et où le retrait est publié. Les DOI sont tous de la forme
\texttt{10.5281/zenodo.} suivi du numéro indiqué.

\section{Les retraits et corrections publiés}

\begin{description}[style=nextline,leftmargin=0pt,font=\normalfont\sffamily\bfseries\small\color{okvermillon},itemsep=5pt]
\item[Juillet 2026 · Les nombres d'un réseau en cosmologie] Le dépôt Agora-K3-DarkMatter publie le 24 juillet une
déclaration qui retire ses nombres « C1/C2 » sur les réseaux. Dans le dépôt des propositions Lean, deux théorèmes
d'« intégralité » se révèlent vides (ils ne disent rien), et les premières annonces sont retirées.
\item[9 septembre · « La première formalisation »] L'article sur l'involution de Fricke (22682139) retire sa
prétention d'être la première formalisation, après qu'un audit a trouvé des travaux antérieurs.
\item[15 septembre · Une interprétation des équations des fluides] L'article « Topological censorship of blow-up »
(22723078) proposait une lecture de la preuve d'OpenAI sur l'explosion des solutions de Navier--Stokes. Trois jours
plus tard, « Where the continuum ends » (22773907) reconnaît les constructions comme des mathématiques valides, retire
l'interprétation et publie un tableau des corrections.
\item[18 septembre · La théorie « dual-scale »] La révision 5 (22837842) rapporte deux échecs : le « verrou » 27720
échoue au test de torsion pour les 25 classes non triviales du groupe M24, ce qui en fait de la numérologie, et la
conjecture testable ne survit pas aux données. La confrontation pré-enregistrée (22837837) rend quatre verdicts
négatifs. Cette ligne de recherche comprenait les articles de juin et juillet qui voulaient expliquer des tensions
dans les observations de galaxies du télescope JWST par des surfaces K3 de la théorie des cordes (20863378, 20863381,
21350629). Les mathématiques Lean de ce programme restent de niveau A et sont séparées de la physique retirée.
\item[20 septembre · Deux « découvertes » sur les fluides quantiques] La bibliothèque QuantumFluids (22853896) retire
deux « découvertes » après une vérification de la littérature.
\item[21 septembre · T-dualité] L'article sur la T-dualité (22872083) retire les affirmations de sa version 1.
\item[21 septembre · Des résultats jamais obtenus] La version 2 de RunuX (22879667) retire les résultats annoncés en
version 1 (sur TPU, avec un modèle de 70 milliards de paramètres, et des certificats Lean), qui n'avaient jamais été
obtenus : « retirés, pas rebaptisés ».
\item[25--27 septembre · L'atelier s'audite] Les audits de l'atelier AutoevolveAI trouvent quatorze composants qui
annonçaient un succès sans avoir fait le travail, dont un script de « relecture » qui n'appelait aucun modèle. L'audit
« affirmations contre preuves » de RunuX (22985926) trouve 141 modules sur 316 qui ne sont que des coquilles vides, et
249 obligations de preuve ouvertes derrière une annonce « zéro sorry ».
\item[3 octobre · Des chiffres de performance] La porte GWAYA (23119953) retire les chiffres de performance de ses
versions 1 et 2.
\item[9 octobre · Le frottement des tourbillons] Le frottement des tourbillons suit la température, pas la densité
normale (23262132) ; l'article sur le transport est corrigé (23262143, version 2.2).
\item[11 octobre · Une coupure de force trop brutale] Dans l'article « par étages » (23293429), la coupure de la force
dans la simulation déplaçait les énergies de plus que l'écart mesuré. L'erratum est publié, et la correction faite
avec la force exacte dans 23646105 (chapitre~\ref{ch:c04}).
\item[19 octobre · Un faux rapport de relecture] Un dépôt (23468458) retire un « rapport de relecture » produit par un
script, avec des relecteurs imaginaires.
\item[21 octobre · « Relu par des pairs »] La seconde édition de l'article de thermodynamique stochastique (23525655)
retire la mention « relu par des pairs » de la première : le fichier de relecture « a été produit par un script qui
n'appelle aucun relecteur ».
\item[21 octobre · Un mécanisme de relecture moléculaire] L'article sur la relecture de Hopfield--Ninio (23526783)
réfute son propre mécanisme « POLE P286R » et ajoute un théorème d'impossibilité.
\item[21 octobre · RK4 « chauffe » les particules] L'édition 2 (23532458) montre que RK4 refroidit au lieu de chauffer,
ce qui réfute la prémisse de l'article sur le stellarator ; ses chiffres de confiance et un théorème sont retirés
(chapitre~\ref{ch:c20}).
\item[21 octobre · Un théorème sur un modèle de turbulence] Le théorème 2.5 de l'article sur le modèle de sous-maille
(23532595) est réfuté, et son affirmation d'explosion retirée.
\item[21 octobre · Un « vérificateur vérifié » de cryptographie] Le « vérificateur SQIsign vérifié » (23532601) accepte
le nombre 256, qu'il aurait dû rejeter ; trois affirmations principales sont retirées.
\end{description}

\section{Nos propres hypothèses réfutées}

Ces résultats ne sont pas des retraits : ce sont des idées du projet, testées comme prévu, et qui ont perdu. Elles sont
publiées comme des résultats.
\begin{itemize}
\item Le théorème de Lieb--Thirring en version matricielle : résultat nul (23245266).
\item Des prédictions réfutées, conservées dans la série sur la conductance (23228685, 23241463, 23244556, 23248393).
\item La série sur les fluides d'excitons : la plupart des tests pré-enregistrés échouent, et les articles le disent.
\item La conjecture QU-1 : réfutée par le calcul pour huit des douze forces testées (23653833,
(chapitre~\ref{ch:c12}), puis par certificat pour deux familles de forces
(23655392, chapitre~\ref{ch:c11}, et 23663551).
\item La formule d'« enveloppe » du projet pour l'hélium : réfutée par les données de diffusion de neutrons, aux
21 points de la région des rotons (23653846, chapitre~\ref{ch:c05}).
\item Une sous-hypothèse pré-enregistrée du théorème isotherme : fausse, et signalée (23660071, chapitre~\ref{ch:c07}).
\item En biologie et en recherche automatique : un certificat dont le transfert aux données de séquences échoue
(23604793) ; des seuils de référence (ceux du « hasard pur ») sous-estimés, corrigés (23610116, 23610123) ; un transfert à la programmation
génétique réfuté (23611425) ; « la recombinaison n'aide pas » (23626924).
\item L'extension en trois dimensions de l'architecture de Sandier--Serfaty : résultat négatif de faisabilité
(23655002) ; une transition quantique attendue : résultat négatif (23655625).
\end{itemize}

\section{Les affirmations excessives restées sans retrait}

Enfin, quelques dépôts du printemps contiennent des affirmations que le dossier ne permet pas de vérifier, et qui n'ont
pas fait l'objet d'un retrait publié. Nous les signalons ici pour que le lecteur ne les prenne pas pour des résultats.
Le même texte, déposé treize fois les 25 et 26 mai (dont 20380024), annonce une boucle d'apprentissage sur TPU, 40\,\%
d'économie d'énergie et une vérification « formelle en Lean~4 » : le dépôt ne contient aucun fichier Lean. Les dépôts
SymBrain (20397720) et CPUGym (20538871, 20539058) rapportent des chiffres de performance qu'un lecteur ne peut pas
vérifier.
